\chapter{The outer and unit-state join for infinitesimal $e$}
\label{component:p05a86be0046c}
\noindent\textbf{Component }\path{p05a86be0046c}\par
{\small\noindent Source: \path{R13/baseline/baseline/baseline/advances/c1_scales/outer_small_e/outer_small_e.tex}\par}
{\small\noindent SHA256: \path{05a86be0046c9b9ee011ff502597dbc9f6c3153d2ee20ef81aa92d23d11a005b}\par}
{\small\noindent Inventory use: SEO1--SEO20 and all exact endpoint/critical-current proofs | SEO1-SEO20, full critical-wall/current identities and endpoints | C1\_\allowbreak{}BARE\_\allowbreak{}OTHER\_\allowbreak{}TERMINALS; precise subordinate uses retained below\par}
\medskip
\section[The outer and unit-state join for infinitesimal e]{The outer and unit-state join for infinitesimal $e$}
\label{p05a86be0046c:sec:small-e-outer}

\begin{lemma}
Use the unchanged physical native C1 source with $c=2$,
$\lambda=n=0$, positive original openings, and nonnegative original
event coefficients. Suppose $k,\zeta,e=k\sqrt{\eta/\zeta}$ are positive
infinitesimals, $\eta$ is a positive bounded unit, and
$r=1-a\eta>0$. Then every independent projective test and state is
tensor SOS except possibly the chart
\[
 u,v,w>0\text{ infinitesimal},\qquad Z=1,
 \qquad w=(W+uY)/(1+u^2).
 \tag{SEO1}
\]
There is no restriction on the positive relative scales of $a,r,J$
or on $k^2/r$. This lemma supplies a join, not a certificate for the
remaining chart (SEO1).
\end{lemma}

\paragraph{The small weak-event scale.}
Put $\kappa=k^2$, $K=1+\kappa$, $H=r+\eta$,
$V_c=\kappa\eta\zeta$, $S=1-\kappa\eta+tV_c$,
$R=1+\kappa r$, $T=K-\kappa\zeta t$,
$M=R-t$, and $J=T-a$. Let $P_4,P_d$ be the short native
polynomials printed in the centered-A outer lemma. All of
$K,H,S,P_4,P_d$ are positive bounded units. Set
\[
 L=KHS,\quad N=KH(1+V_c)-J\eta P_d,\quad
 d_0=4\eta P_d,\quad A_0=2\eta^2P_4,\quad D_0=2N.
\]
After exact cancellation of the native $r$ factors, the source is
\[
 d=\frac{ad_0}{L},\qquad A=\frac{aA_0}{rL},\qquad
 D=\frac{aD_0}{JL},\qquad
 \frac{m}{ae^2}=
 \frac{2(1-t\zeta)(1+V_c)}{MS}.                 \tag{SEO2}
\]
Here $\bar N=1$. Native $\beta\ge0$ gives the exact bound
\[
 \frac{t}{e^2}\le\frac{2c_b}{\eta c_d}
 =\frac{K[1+(1-t)\zeta]H}{\eta P_d}.
\]
Thus $t$ is infinitesimal, $M\sim1$, and $m/(ae^2)\sim2$.
In particular
\[
 \frac{k^2}{e^2}=\frac\zeta\eta
 \quad\hbox{is infinitesimal}.                       \tag{SEO3}
\]
The ratios $d/a,A/a,D/a$ have positive real lower bounds, while
$d/a$ is a bounded unit. These statements allow arbitrary positive
opening rates and literal $\beta=0$.

\paragraph{Choose the critical parameter without a small denominator.}
There are two source cases. If $a$ is a bounded unit, keep $A,D,d$
and replace $c=2$ by the critical number
\[
 c_* =\frac{d(d+2A)(d+2D)}{2(4AD-d^2)}.
 \tag{SEO4}
\]
This is positive and bounded, and $c_*-2=O(k^2)$.
Here is an exact short verification of the potentially delicate bound.
Define
\[
 Q=4A_0D_0-rJd_0^2,
\]
so $\bar Q=16\bar\eta^2\bar a\bar H>0$. Then
\[
 c_*-2=
 \frac{ad_0(d_0r+2A_0)(d_0J+2D_0)-4QL}{2QL}.
 \tag{SEO5}
\]
After $r=1-a\eta$ is inserted, its polynomial numerator vanishes
identically at $\kappa=0$, with $\zeta$ still free. It is therefore
$\kappa$ times a polynomial in bounded native variables; its
denominator is a positive unit. No inverse of $r$ or $J$ occurs.

If $a$ is infinitesimal, then $r,J\sim1$. Divide the whole parent
by $a$ and replace the original $A$ by
\[
 A_* =\frac{d^2(4+d+2D)}{2(8D-2dD-d^2)}.       \tag{SEO6}
\]
The denominator is positive. In this case
$A_*/a,D/a,d/a$ are positive bounded units, $c/a=2/a$ is
unbounded, and $(A_*-A)/a=O(k^2)$. To verify this without a
rate comparison, put
\[
 B_0=8D_0L-a(2d_0D_0+d_0^2J).
\]
Its residue is $16\bar H>0$, and
\[
 \frac{A_*-A}{a}=
 \frac{r d_0^2\{4JL+a(d_0J+2D_0)\}-2A_0B_0}
 {2rLB_0}.                                           \tag{SEO7}
\]
Again the numerator vanishes identically at $\kappa=0$, so it has
an exact polynomial factor $\kappa$, with a unit denominator.

Use $\sigma=1$ in the first case and $\sigma=a$ in the second,
and henceforth normalize all parent coefficients by $\sigma$.
Let $\epsilon$ be the positive part of the critical increase in the
chosen normalized coefficient. From (SEO2)--(SEO7),
\[
 \epsilon\text{ is infinitesimal},\qquad
 \frac{\epsilon}{m/\sigma}=O(\zeta)\text{ is infinitesimal}.
 \tag{SEO8}
\]
If the chosen critical value is smaller than the actual coefficient,
retain the positive excess packet instead. In either source case
the chosen critical parent obeys
\[
 2c(4AD-d^2)=d(d+2A)(d+2D),                  \tag{SEO9}
\]
with the letters in (SEO9) now denoting the normalized critical
coefficients. In the first case $c,d$ are units and $A,D$ are
bounded below by positive units; in the second $A,D,d$ are units
and $c$ may be unbounded.

\paragraph{A critical-wall frame permitting the large coefficients.}
Write $D_u=1+u^2$, $D_v=1+v^2$ and
\[
 B=X+uZ,\quad U=uW,\quad V=vZ,\quad J_t=U-Y,
 \quad C=4D+2d,
\]
and define the pure-coefficient numbers
\[
 y=\frac d{2A},\quad
 b=\frac{4AD-d^2}{AC},\quad s=1-b,\quad q=y+b.
\]
The critical identity gives the useful exact relations
\[
 b=\frac C{4c+C},\quad Cs=4cb,\quad
 cb=\frac{d(2A+d)}{4A}.                         \tag{SEO10}
\]
In both source cases, $y,b,s,q$ are bounded and $b,s>0$.
The entire signed core plus the three critical packets is exactly
\begin{align*}
 F_*={}&4cD_vJ_t^2+4AD_uD_vY^2
       +C(J_t+uvB-D_uV)^2+2dD_vB^2\\
       &+2dD_u\{v^2Y^2-2YV-V^2\}.              \tag{SEO11}
\end{align*}

In the first source case use
\[
 H_A=\sqrt A(Y-yV),\qquad
 H_D=\sqrt C(J_t+uvB-(b+u^2)V).
 \tag{SEO12}
\]
The inverse is $Y=yV+H_A/\sqrt A$ and
$J_t=(b+u^2)V-uvB+H_D/\sqrt C$.
All coefficients in (SEO11) after this substitution are bounded.
Indeed the large $A,C$ occur only in the normalized squares and in
the bounded combinations $Ay^2$, $\sqrt A\,y$, $Cs^2$,
$\sqrt C\,s$; use $Cs=4cb$ and bounded $c$ for the last two.
At $u=v=0$, the hard block on $(B,H_A,H_D)$ is
\[
 \operatorname{diag}(2d,4,1+4c/C)\succ0.
\]
On its exact zero section the scalar coefficient of $V^2$ is
\[
 g u^2+\gamma v^2+4cu^4+g(y+2)u^2v^2+4cu^4v^2,
 \quad g=\frac{d(2A+d)}A,\quad
 \gamma=\frac{Dd(2A+d)^2}{A^2(2D+d)}.             \tag{SEO13}
\]
Both $g,\gamma$ are positive bounded units. Every hard-$V$ entry
has angular degree at least two, with bounded coefficients; this is
an exact full-parent polynomial statement. Allocate it on $uV,vV$,
leaving a factor $u$ or $v$ in its coefficient. The resulting exact
five-root Gram on $(B,H_A,H_D,uV,vV)$ therefore has strict residue.

In the second source case set $L_c=4c+C$ and instead use
\[
 H_A=\sqrt A(Y-yV),\qquad
 H_c=\sqrt{L_c}(J_t-bV).                         \tag{SEO14}
\]
Its coefficients are bounded because $A,D,d$ are units and
$c/L_c$, $cb/\sqrt{L_c}$ and $cb^2$ are bounded by (SEO10).
Its hard block is $\operatorname{diag}(2d,4,1)$.
The exact zero-section scalar is
\[
 g u^2+\gamma v^2+C u^4+gy\,u^2v^2.
 \tag{SEO15}
\]
The same complete hard-$V$ allocation gives a strict five-root
Gram on $(B,H_A,H_c,uV,vV)$. Thus this is a full critical-wall
certificate in both source cases, even with the stated unbounded
coefficients. Every radical and denominator used in the roots is
pure coefficient; no mixed inverse is used.

\paragraph{All unit-$W$ states at the small-angle corner.}
Normalize $W=1$, with $Z,Y,X$ bounded, and take $u,v$ infinitesimal.
Retain the complete original $\alpha,\beta$ packets. Put
\[
 R_m=v-X+uvY-uZ=v(1+uY)-B,\qquad
 V=\frac{Z}{1+uY}(R_m+B).                        \tag{SEO16}
\]
The displayed denominator is a positive pure-testing unit.
In the first source case the only negative correction to the
critical parent is $-4\epsilon D_vJ_t^2$; (SEO12) and (SEO16)
express $\sqrt\epsilon J_t$ as a bounded combination of the five
critical roots and $\sqrt{m/\sigma}\,R_m$, with every coefficient
having zero residue. In the second case the only negative
correction is $-4\epsilon D_uD_vY^2$; use
$Y=yV+H_A/\sqrt A$ and (SEO16) for the identical conclusion.
The only possibly delicate coefficient is a bounded testing
multiple of $\sqrt{\epsilon/(m/\sigma)}$, and it is infinitesimal
by (SEO8). The extra original $m$ square contributes the strict
diagonal $4$. Hence the complete perturbed Gram remains strict.
This includes unit or infinitesimal $Z$, literal $Z=0$, and both
signs of the bounded testing state. It uses the entire original
weak $m$ event, without a false positive-unit lower bound on $m$.

\paragraph{Finite outer angles without the weak event.}
The same critical parent also gives the two finite outer charts.
In (SEO11) regard $B,Y,J_t,V$ as independent roots. Exactly,
\[
 F_*(0,v)=F_*(0,0)
       +v^2\{2dB^2+(4A+2d)Y^2+4cJ_t^2\}.       \tag{SEO17}
\]
The zero-angle form has three positive hard pivots and null ray
$(B,Y,J_t,V)=(0,y,b,1)V$. On this null ray the additional
coefficient is $\gamma v^2>0$. The normalized hard frames
(SEO12) or (SEO14) keep every coefficient bounded. Thus when
$u$ is infinitesimal and $v$ a bounded unit the full four-root
residue is strict. The infinitesimal critical-parameter deficit
is absorbed directly on these four roots; retain the entire
$m,\alpha,\beta$ packets.

For bounded-unit $u$ and infinitesimal $v$, set $V'=D_uV$.
The other exact identity is
\[
 F_*(u,0)=F_*(0,0;V')
        +u^2\{4AY^2+2d(V')^2/D_u\}.             \tag{SEO18}
\]
Use the same normalized critical roots with $V'$ in place of $V$.
In the first source case include $uvB$ in the $H_D$ root so that
the unbounded $D$ event remains isolated for nonzero small $v$.
The extra null mass in (SEO18) is positive in residue because
$d$ and $u$ are units. The bounded full four-root Gram is strict,
and the deficit is again a zero-residue correction. These two
identities include every state, not merely the normal restriction.

For completeness the same argument covers both finite angles of
nonzero residue, without a separate physical-source admission.
The exact determinant of (SEO11) on $(B,Y,J_t,V)$ is
\[
 64cdD_u^2D_v^2
 \{A(2D+d)u^2+D(2A+d)v^2\}.                    \tag{SEO20}
\]
The determinant in the normalized frame (SEO12) is this number
divided by $AC$, namely
\[
 32cdD_u^2D_v^2
 \left\{u^2+\frac{D(2A+d)}{A(2D+d)}v^2\right\}.
\]
It is a positive unit whenever at least one finite angle has
nonzero residue. The hard restriction $V=0$ is explicitly
\[
 \left(4D_uD_v+\frac{2dD_uv^2}{A}\right)H_A^2
 +H_D^2+4cD_v(H_D/\sqrt C-uvB)^2+2dD_vB^2,
\]
which has a strict bounded residue. Hence its full residue Gram
is positive definite by (SEO20). In frame (SEO14), the determinant
is (SEO20) divided by $A(4c+C)$. Here $A,D,d$ are units and
$c/(4c+C)\sim1/4$, so it too is a positive unit. Its hard
restriction is
\[
 \left(4D_uD_v+\frac{2dD_uv^2}{A}\right)H_A^2
 +\frac{4cD_v}{4c+C}H_c^2
 +C(H_c/\sqrt{4c+C}+uvB)^2+2dD_vB^2,
\]
again strict. All entries of both full transformed Grams are
bounded as already proved. The infinitesimal deficit is a
zero-residue correction on these four roots. This settles every
finite angle pair outside the double-infinitesimal corner.

\paragraph{Projective infinity charts.}
For these charts normalize by $a$, clip the positive coefficients
$c/a,A/a,D/a$ at suitable positive real lower bounds, and retain
their positive original excess squares. Retain the whole original
$m,\alpha,\beta$ packets. The normalized $d/a$ is a unit.
All the infinity-infinity and infinity-zero four-root identities
in (BE2) and the subsequent outer routing apply without $m$.
For $u=\infty$ in residue and $v$ a bounded unit, the signed-core
residue itself retains
$2d\{(W+vX)^2+(vY-Z)^2\}$; together with $cW^2,AY^2$ it is
strict, so no $m(vY-Z)^2$ lower bound is needed.
For $v=\infty$ and $u$ a bounded unit the old signed-core and
$c,D$ root argument is unchanged.
At the remaining $(u,1/v)=(0,0)$ corner, replace the obsolete
state-$W$ mass argument by the exact root ideal
$(X,Y,U=uW,Z)^2$ after retaining $m$. In $q^2F(u,1/q)$ its
complete selected residue is
\[
 2dX^2+4c(U-Y)^2+(4A+2d)Y^2+4DZ^2\succ0.
 \tag{SEO19}
\]
Every remaining occurrence of $W$ carries $u$, so this is a full
bounded Gram including all states. Literal infinity is included
by these homogeneous identities. Literal zero axes have their
unchanged original-event SOS. Finite tests with both noncommuting
residues were included directly in (SEO20).

\paragraph{The residual state chart.}
Only $u,v$ infinitesimal remain. Normalize the state by a maximal
coordinate. On a negative pure-testing-core branch, $X,Y$ have
zero residue. The unit-$W$ chart was just proved for every bounded
$Z$. Otherwise $Z$ is a unit, so normalize $Z=1$ and $W$ is
infinitesimal. Then $w=(W+uY)/D_u$ is infinitesimal. The inherited
exact testing Schur identity and sign gauges give (SEO1).
Complex states split into their real and imaginary evaluations.
The exact checker accompanying this lemma verifies its source
cancellations, critical roots, full hard-normal orders, both finite
outer identities, and the repaired projective-corner identity.
